\documentclass[]{article}
\usepackage[utf8]{inputenc}
\usepackage{geometry}
\usepackage{enumitem}
\usepackage{amsmath}
\usepackage{amssymb}
\usepackage{amsfonts}
\usepackage{makeidx}
\usepackage{booktabs}
\usepackage{tabularx}
\usepackage{caption}
\usepackage{textcomp}
\usepackage{authblk}
\usepackage{graphicx}
\usepackage{caption}
\usepackage{subcaption}
\usepackage{multicol}
\usepackage{multirow}
\usepackage{array}
% \usepackage{fourier} 
\usepackage{makecell}
\usepackage{fancyhdr}
%\usepackage{showframe}
% \usepackage{fancyheadings}
	
%\pagestyle{fancy}
\fancyhf{}
\lhead{GPU Technology Conference, March 26-29, 2018, Silicon Valley, USA}
%\rhead{\textbf{Accepted for presentation}}
\renewcommand{\headrulewidth}{0.4pt}
\newcommand{\RN}[1]{%
	\textup{\uppercase\expandafter{\romannumeral#1}}%
}
\renewcommand{\tabularxcolumn}{m}
\renewcommand\theadalign{bc}
\renewcommand\theadfont{\bfseries}
% \renewcommand\theadgape{\Gape[4pt]}
% \renewcommand\cellgape{\Gape[4pt]}
%%%%%%%%%%%%%%%%%%%%%%%%%%%%%%%%%%%%%%%%%%%%%%%%%%%%%%%%%%%%%%%%%%%%%%%%%%%%%
\geometry{	a4paper, total={170mm,257mm}, left=20mm, top=20mm, includehead}
%\geometry{	a4paper, total={170mm,257mm}, left=20mm, top=20mm}
%\geometry{a4paper,left=3cm,right=2cm,top=2.5cm,bottom=2.5cm}
%\geometry{a4paper, margin=2cm, includehead}
%%%%%%%%%%%%%%%%%%%%%%%%%%%%%%%%%%%%%%%%%%%%%%%%%%%%%%%%%%%%%%%%%%%%%%%%%%%%%%%%%%%%%%%%%%%%%%%%%%%%%%%%%%%%%
%opening
\title{\textbf{GPU-acceleration of an unstructured compressible DNS code for realistic applications}}
\author{N. H. Maruthi, Rajesh Ranjan, S. M. Deshpande\thanks{Corresponding author}, Roddam Narasimha}
\affil{
  Engineering Mechanics Unit\\
  Jawaharlal Nehru Centre for Advanced Scientific Research\\
  Bengaluru, 560064\\
  Email: smd@jncasr.ac.in
}
\author{Bharatkumar Sharma, Shankara Rao Thejaswi Naditale}
        \affil{NVIDIA, Bengaluru, 560045 \\
       Email: snanditale@nvidia.com
}

\date{}
\begin{document}
\maketitle
\thispagestyle{fancy}
%%%%%%%%%%%%%%%%%%%%%%%%%%%%%%%%%%%%%%%%%%%%%%%%%%%%%%%%%%%%%%%%%%%%%%%%%%%%%%%%%%%%%%%%%%%%%%%%%%%%%%%%%%%%%%%%%%%%%
\renewcommand{\abstractname}{Session description}
\begin{abstract}
 In this work, we present the methodologies and optimization techniques used in converting an in-house compressible Direct Numerical Simulation (DNS) code from C++ to GPU-enabled CUDA. This code, named ANUROOP, has been developed to simulate the complex flow past gas turbine blades in an unstructured finite-volume framework. The performance results of the code on modern CPUs and GPUs is reported here. These performance results show that when the code is optimally accelerated, ANUROOP on a single P100 is 110x faster compared to Intel’s Xeon E5-2680 v3 (Haswell) single core, and 3.1x faster compared to one full Haswell node (24 cores). Thus GPUs provide remarkable speed-up without any loss of accuracy. Therefore GPUs can be used for computing complex fluid-dynamical problems to get accurate results in a realistic time-frame.
\end{abstract}
%%%%%%%%%%%%%%%%%%%%%%%%%%%%%%%%%%%%%%%%%%%%%%%%%%%%%%%%%%%%%%%%%%%%%%%%%%%%%%%%%%%%%%%%%%%%%%%%%%%%%%%%%%%%%%%%%%%%%
%%%%%%%%%%%%%%%%%%%%%%%%%%%%%%%%%%%%%%%%%%%%%%%%%%%%%%%%%%%%%%%%%%%%%%%%%%%%%%%%%%%%%%%%%%%%%%%%%%%%%%%%%%%%%
\section*{Extended abstract}
In this work, we present the methodologies and optimization techniques used in converting an in-house compressible Direct Numerical Simulation (DNS) code from C++ to GPU-enabled CUDA. This code, named ANUROOP, has been developed at JNCASR, Bengaluru to simulate the complex flow past gas turbine blades in an unstructured finite-volume framework. 
To assess the performance of the GPU-accelerated ANUROOP, we have conducted weak and strong scaling experiments. Weak scaling experiments are performed on an IBM Power System S822LC cluster code-named Minsky (USA), and strong scaling experiments are conducted on both Minsky and DHRUVA-3 (DRDO ANURAG Lab, India) clusters.  Minsky has four P100 GPUs in each node with a total of 34 nodes, while DHRUVA-3 has two K20 GPUs in each node with a total 114 nodes. The CPU version of ANUROOP was benchmarked on SahasraT (CRAY XC40, IISc, Bengaluru), which has Intel Xeon E5-2680 v3 (Haswell) processors with 24 cores in each node and has a total of 1376 nodes. 
For all the benchmark studies reported here, flow past the widely studied low-pressure turbine blade T106A has been chosen as the test case, and the runtime of the first 1000 iterations (excluding I/O time) is taken for the performance analysis.
It is demonstrated that on a single P100 GPU, ANUROOP is approximately 110x faster compared to a run on a single-CPU on SahasraT, while when compared to one full node (24 CPU cores), the speed-up obtained is 3.1x. In the comparison between GPUs, we have observed that the code on the P100 runs 2.7x faster compared to a K40 GPU. 
%%%%%%%%%%%%%%%%%%%%%%%%%%%%%%%%%%%%%%%%%%%%%%%%%%%%%%%%%%%%%%%%%%%%%%%%%%%%%%%%%%%%%%%%%%%%%%%%%%%%%%%%%%%%%
%%%%%%%%%%%%%%%%%%%%%%%%%%%%%%%%%%%%%%%%%%%%%%%%%%%%%%%%%%%%%%%%%%%%%%%%%%%%%%%%%%%%%%%%%%%%%%%%%%%%%%%%%%%%%
\end{document}
